\documentclass[10pt,a4paper]{amsart}
\usepackage[margin=19mm]{geometry}
\usepackage{amsmath,amssymb,mathtools}
\usepackage[scr=rsfs,cal=euler]{mathalfa}
\usepackage{xfrac}
\usepackage[unicode,hidelinks,bookmarks=false]{hyperref}
\newcommand{\ie}{i.e.\ }
\newcommand{\cf}{cf.\ }
\newcommand{\resp}{respectively\ }
\newcommand{\Subsection}{\subsection}
\providecommand{\nobreakdash}{\nobreak\hbox{-}}
\setlength{\emergencystretch}{2em}
\multlinegap0pt
\usepackage{bm}
\usepackage{bbm}
\usepackage{dsfont}
\usepackage{xspace}
\usepackage{cancel}
\usepackage{mathtools}
\usepackage{amsthm}
\makeatletter
\@ifundefined{theorem}{\newtheorem{theorem}{Theorem}}{}
\@ifundefined{proposition}{\newtheorem{proposition}[theorem]{Proposition}}{}
\makeatother
\newcommand{\dvh}[0]{\mathrm{div}_{h}\,}
\newcommand{\nablah}[0]{\nabla_{h}}
\title[Local well-posedness of strong solutions to the non-isentrop]{Local well-posedness of strong solutions to the non-isentropic compressible primitive equations with vertical diffusion}
\author{Rupert Klein, Jinkai Li, Xin Liu, Edriss S. Titi}
\date{Source: arXiv:2602.19801v1 (CC BY 4.0)}
\begin{document}
\maketitle

\noindent\emph{Source and licence.} Local well-posedness of strong solutions to the non-isentropic compressible primitive equations with vertical diffusion. Version arXiv:2602.19801v1 (CC BY 4.0), distributed under the Creative Commons Attribution 4.0 International licence, as recorded in the arXiv licence metadata for this exact version and shown on the abstract page https://arxiv.org/abs/2602.19801v1. Full rights notice: licenses/arxiv-2602.19801-CC-BY-4.0.txt.\par\smallskip

\noindent\textit{Editorial scope.} Counted unit: the Proposition whose source label is \texttt{PROP-EGYINEQT} in the article (source0.tex lines 561--575, source label \texttt{PROP-EGYINEQT}), delivered together with the article's own complete original proof of it (source0.tex lines 577--738, one \texttt{proof} environment of 162 lines). No mathematics is written, repaired or re-derived: statement and proof are the article's own text line for line. Required context retained uncounted: Q (lines 150--154); phi (lines 221--223); BC1 (lines 252--255); BC2 (lines 252--255); IC (lines 257--260); COME (lines 331--333); EQp (lines 369--375); EQs (lines 369--375); EQv (lines 369--375); EQw (lines 369--375); ASSUMIC1 (lines 379--383); ASSUMIC2 (lines 379--383); PROPLOC-e (lines 388--401); Phi1 (lines 410--414); Phi2 (lines 410--414); Phi3 (lines 410--414); EQs' (lines 417--422); EQv' (lines 417--422); PROP-ESTPhis (lines 426--449); ASSUM-APRM (lines 431--435); ASSUM-LBsp (lines 431--435). Every definition, display and label of those lines is retained unchanged. Omitted from this package and not redistributed: 1--149, 155--220, 224--251, 256--256, 261--330, 334--368, 376--378, 384--387, 402--409, 415--416, 423--425, 436--560, 576--576, 739--2287. Nothing inside a retained excerpt is omitted or altered: every retained body is delivered exactly as the article's lines are written, so no omission receipt of any kind accompanies this package. Citations: the retained excerpts carry no citation command at all, so no bibliography entry of the article is transcribed and no reference is redistributed. Reference adaptation: none was needed and none was made; every reference inside the delivered unit resolves inside it, so no unresolved reference or citation is produced. Printed slips kept literal: no printed word, symbol, hypothesis or number of the counted statement or of its proof was altered. Layout and class: the article's own class is \texttt{amsart} with packages including mathrsfs, amssymb, amsmath, amsfonts, amsthm, natbib, graphicx, bm, xy, array, color; the wrapper is the lane's pinned \texttt{amsart} editorial wrapper at 10pt on A4, which restates the theorem-like environments the retained text needs and adds the article's own macro definitions that the retained text needs (\texttt{\textbackslash dvh}, \texttt{\textbackslash nablah}), with extra wrapper packages (bm, bbm, dsfont, xspace, cancel, mathtools). The delivered numbering is the unit's own. Assets: no retained span contains a figure environment, a graphics-inclusion command, a PDF-page inclusion, a table or a TikZ picture; no image file is copied into this package, the package declares no asset dependency and no image file is redistributed with it. Licence: version arXiv:2602.19801v1 of this article is distributed under the Creative Commons Attribution 4.0 International licence, as recorded in the arXiv licence metadata for this exact version and shown on the abstract page https://arxiv.org/abs/2602.19801v1. A full rights notice is delivered at licenses/arxiv-2602.19801-CC-BY-4.0.txt. Characters: every retained excerpt is pure ASCII, so the delivered engine, XeLaTeX with its default Latin Modern text font, reports no missing character.
\begin{eqnarray}
	&\mathbb S_h = \mu (\nablah v + \nablah^\top v) + \lambda \dvh v \mathbb I_2, \quad
	p = R \rho \theta, \nonumber\\
  &Q(\nabla v)=\mathbb S_h:\nabla_hv+\mu|\partial_zv|^2,\label{Q}
\end{eqnarray}

\begin{eqnarray}
  \phi(v,p):=\text{div}_h\tilde v+\frac{1}{\gamma p}\left(\tilde v\cdot\nabla_hp-(\gamma-1)\widetilde{Q(\nabla v)}\right).\label{phi}
\end{eqnarray}

\begin{eqnarray}
  &v, \sigma, \mbox{ and } p \mbox{ are periodic in }x, y, \label{BC1}\\
  &\partial_zv|_{z=0,1}=0, \quad\partial_z\sigma|_{z=0,1}=0. \label{BC2}
\end{eqnarray}

\begin{equation}
  \label{IC}
  (v, \sigma, p)|_{t=0}=(v_{0}, \sigma_{0}, p_{0}).
\end{equation}

\begin{equation}
  \|[D^\alpha, f]g\|_q\leq C\left(\|\nabla f\|_{r_1}\|g\|_{W^{m-1, s_1}}+\|g\|_{r_2}\|f\|_{W^{m,s_2}}\right),\label{COME}
\end{equation}

\begin{eqnarray}
  &\partial_tv+(v\cdot\nabla_h)v+w\partial_zv+\sigma\nabla_hp=\mu\sigma\Delta v+(\mu+\lambda)\sigma\nabla_h\text{div}_hv,\label{EQv}\\
  &w=\nu\partial_zv-\int_0^z\phi(v,p)dz',\label{EQw}\\
  &\partial_t\sigma+v\cdot\nabla_h\sigma+w\partial_z\sigma+\sigma(\phi(v,p)-\text{div}_hv)=\nu\sigma\partial_z^2\sigma+\epsilon\Delta_h\sigma,
  \label{EQs}\\
  &\partial_tp+\bar v\cdot\nabla_hp+\gamma p\text{div}_h\bar v=\epsilon\Delta_hp+(\gamma-1)\overline{Q(\nabla v)}, \label{EQp}
\end{eqnarray}

\begin{eqnarray}
    (v_{0}, \sigma_{0})\in H^3(\mathcal O),\quad p_{0}\in H^3(\mathbb T^2),\quad
    \sigma_{0}\geq\underline\sigma, \quad p_{0}\geq\underline p,\label{ASSUMIC1}\\
    \partial_zv_{0}|_{z=0,1}=0,\quad \partial_z\sigma_{0}|_{z=0,1}=0.\label{ASSUMIC2}
\end{eqnarray}

\begin{proposition}
  \label{PROPLOC-e}
  Given $(v_{0}, \sigma_{0}, p_{0})$ satisfying (\ref{ASSUMIC1})--(\ref{ASSUMIC2}).
  Then, there is a positive time $T_\epsilon$,
  depending only on $\epsilon$, $\underline\sigma$, $\underline p$, and
  $\|(v_{0},\sigma_{0})\|_{H^3(\mathcal O)}+\|p_{0}\|_{H^3(\mathbb T^2)}$, such that system (\ref{EQv})--(\ref{EQp}),
  subject to (\ref{BC1})--(\ref{IC}), has a unique local in time solution $(v,\sigma,p)$ on $\mathcal O\times(0,T_\epsilon)$,
  satisfying
  \begin{eqnarray*}
  &&\inf_{(x,y,z,t)\in\mathcal O\times(0,T_\epsilon)}\sigma>0,\quad \inf_{(x,y,t)\in\mathbb T^2\times(0,T_\epsilon)}p>0,\\
    &&(v,\sigma)\in C([0,T_\epsilon]; H^3(\mathcal O))\cap L^2(0,T_\epsilon; H^4(\mathcal O)),\quad(\partial_tv,\partial_t\sigma)\in L^2(0,T_\epsilon; H^2(\mathcal O)),\\
    &&p\in C([0,T_\epsilon]; H^3(\mathbb T^2))\cap L^2(0,T_\epsilon; H^4(\mathbb T^2)),\quad\partial_tp\in  L^2(0,T_\epsilon; H^2(\mathbb T^2)).
  \end{eqnarray*}
\end{proposition}

\begin{eqnarray}
  &&\Phi_1=\Phi_1(v,\sigma,p):=-(v\cdot\nabla_h)v-w\partial_zv-\sigma\nabla_hp,\label{Phi1}\\
  &&\Phi_2=\Phi_2(v,\sigma,p):=-w\partial_z\sigma+\sigma(\text{div}_hv-\phi(v,p)),\label{Phi2}\\
  &&\Phi_3=\Phi_3(v,p):=(\gamma-1)\overline{Q(\nabla v)}-\gamma p\text{div}_h\bar v, \label{Phi3}
\end{eqnarray}

\begin{eqnarray}
  \partial_tv-\mu\sigma\Delta v-(\mu+\lambda)\sigma\nabla_h\text{div}_hv&=&\Phi_1,\label{EQv'}\\
  \partial_t\sigma+v\cdot\nabla_h\sigma-\nu\sigma\partial_z^2\sigma-\epsilon\Delta_h\sigma&=&\Phi_2,
  \label{EQs'}\\
  \partial_tp+\bar v\cdot\nabla_hp-\epsilon\Delta_hp&=&\Phi_3. \label{EQp'}
\end{eqnarray}

\begin{proposition}
\label{PROP-ESTPhis}
%\marginpar{\color{blue}Jinkai: How do you want to define `smooth solution'?}
Given a positive time $\mathcal T$ and let $(v,\sigma, p)$ be a solution to system (\ref{EQv})--(\ref{EQp}), subject to (\ref{BC1})--(\ref{IC}), on $\mathcal O\times(0,\mathcal T)$, enjoying the regularities
stated in Proposition \ref{PROPLOC-e}. Assume that
\begin{eqnarray}
  \sup_{0\leq t\leq\mathcal T}\left(\|v\|_{H^3(\mathcal O)}^2+\|\sigma\|_{H^2(\mathcal O)}^2+\|p\|_{H^3(\mathbb T^2)}^2\right)(t)\leq M,  \label{ASSUM-APRM}\\
  \inf_{(x,y,z,t)\in\mathcal O\times(0,\mathcal T)}\sigma\geq 0.5\underline\sigma,\quad \inf_{(x,y,t)\in\mathcal T^2\times(0,\mathcal T)}
  p\geq 0.5\underline p,\label{ASSUM-LBsp}
\end{eqnarray}
for some positive number $M$.
Let $\Phi_1, \Phi_2,$ and $\Phi_3$ be given by (\ref{Phi1}),
(\ref{Phi2}), and (\ref{Phi3}), respectively.
Then, the following estimates hold
\begin{align*}
  &\|(\Phi_1, \Phi_2, \Phi_3)\|_1\leq C, \quad \|\Phi_1\|_{H^2}\leq \nu C_*\|\partial_zv\|_{H^2}\|\partial_z\sigma\|_{H^2}+C,
  \\
  &\|\Phi_2\|_{H^2}\leq C\left(1+\|\partial_z\sigma\|_{H^2}^\frac32\right), \quad
  \|\Phi_3\|_{H^3}\leq C(1+\|v\|_{H^4}),
\end{align*}
for any $t\in(0,\mathcal T)$,
where $C_*$ is an absolute positive constant, and $C$ is a positive constant depending only on $\gamma$, $\nu$, $\mu$,
$\lambda$, $\underline p$, and $M$.
\end{proposition}

\begin{proposition}
\label{PROP-EGYINEQT}
Assume that all the assumptions stated in Proposition \ref{PROP-ESTPhis} hold. Then, it holds that
\begin{align*}
  &\frac{d}{dt}\|(v,\nabla\Delta v)\|_2^2+\mu\underline\sigma\|\Delta^2v\|_2^2
  \leq \nu C_*\|\partial_z\sigma\|_{H^2}\|v\|_{H^4}\|\partial_zv\|_{H^2}+C\left(1+\|v\|_{H^4}^\frac32\right),\\
  &\frac{d}{dt}\|(\sigma,\Delta\sigma)\|_2^2+\nu\underline\sigma\|\Delta\partial_z\sigma\|_2^2
  +\epsilon\|(\nabla_h\sigma,\nabla_h\Delta\sigma)\|_2^2
  \leq C\left(1+\|\partial_z\sigma\|_{H^2}^\frac32\right),\\
  &\frac{d}{dt}\|(p,\nabla_h\Delta_hp)\|_2^2+\epsilon\|(\nabla_hp,\nabla_h^2\Delta_hp)\|_2^2\leq C\left(1+\|v\|_{H^4}\right),
\end{align*}
for any $t\in(0,\mathcal T)$,
where $C_*$ is an absolute positive constant, and $C$ is a positive constant depending only on $\gamma$, $\nu$, $\mu$,
$\lambda$, $\underline p$, and $M$.
\end{proposition}

\begin{proof}
Noticing that $\int_0^1\tilde fdz=0$ and since $p$ is independent of $z$, it is clear from the expression of $\phi(v,p)$ in (\ref{phi}) that
$\int_0^1\phi(v,p)dz=0.$
Thus, recalling (\ref{EQw}) and by the boundary condition (\ref{BC2}), one has
\begin{equation}
  \label{BCw}
  w|_{z=0,1}=0.
\end{equation}
Besides, recalling (\ref{Q}) and using the boundary condition $\partial_zv|_{z=0,1}=0$, it holds that
\begin{equation}
  \label{BCphi2}
  \partial_z\phi(v,p)|_{z=0,1}=0.
\end{equation}
Taking $\partial_z$ to \eqref{EQv} and \eqref{EQs}, and recalling (\ref{BC2}), \eqref{BCw}, and \eqref{BCphi2}, one gets
\begin{equation}
  \label{BC3}
  \partial_z^3v|_{z=0,1}=0,\quad\partial_z^3\sigma|_{z=0,1}=0.
\end{equation}

We first carry out the energy inequalities for $v$. Multiplying \eqref{EQv'} with $v$, integrating the resultant over $\mathcal O$, one obtains by Proposition \ref{PROP-ESTPhis} that
\begin{eqnarray}
  \frac12\frac{d}{dt}\|v\|_2^2&=&\int\sigma(\mu\Delta v+(\mu+\lambda)\nabla_h\text{div}_hv)\cdot vdxdydz+\int\Phi_1\cdot vdxdydz \nonumber\\
  &\leq& C\|\sigma\|_2\|\nabla^2v\|_2\|v\|_\infty+\|\Phi_1\|_1\|v\|_\infty\nonumber\\
  &\leq& C\|\sigma\|_{2}\|v\|_{H^2}^2+\|v\|_{H^2}\leq C.\label{EQ16}
\end{eqnarray}
Applying $-\Delta$ to \eqref{EQv'}, multiplying the resultant with $\Delta^2v$, integrating over $\mathcal O$, and using the boundary
conditions \eqref{BC1}, \eqref{BC2}, and \eqref{BC3}, it follows from integrating by parts that
\begin{eqnarray}
  \frac12\frac{d}{dt}\|\nabla\Delta v\|_2^2
  &=&-\int\Delta(\mu\sigma\Delta v+(\mu+\lambda)\sigma\nabla_h\text{div}_hv+\Phi_1)\cdot\Delta^2vdxdydz\nonumber\\
  &=&-\int(\mu\sigma\Delta^2v+(\mu+\lambda)\sigma\nabla_h\text{div}_h\Delta v+\Delta\Phi_1)\cdot\Delta^2vdxdydz\nonumber\\
  &&-\int(\mu[\Delta,\sigma]\Delta v+(\mu+\lambda)[\Delta,\sigma]\nabla_h\text{div}_hv)\cdot\Delta^2v dxdydz. \label{EQ17}
\end{eqnarray}
Using the boundary conditions (\ref{BC1}), \eqref{BC2}, and \eqref{BC3}, it follows from integrating by parts that
\begin{eqnarray*}
  -\int\sigma\nabla_h\text{div}_h\Delta v\cdot\Delta^2vdxdydz
  =\int(\sigma\text{div}_h\Delta v\Delta^2\text{div}_hv+\text{div}_h\Delta v\nabla_h\sigma\cdot\Delta^2v)dxdydz\\
  =-\int[\sigma|\nabla\text{div}_h\Delta v|^2+\text{div}_h\Delta v(\nabla\sigma\cdot\nabla\Delta\text{div}_hv+\nabla_h\sigma\cdot\Delta^2v)]dxdydz.
\end{eqnarray*}
Substituting this into (\ref{EQ17}), recalling \eqref{ASSUM-LBsp}, applying Proposition \ref{PROP-ESTPhis}, and noticing that $\mu+\lambda>0$,
one deduces
\begin{eqnarray}
  &&\frac12\frac{d}{dt}\|\nabla\Delta v\|_2^2+0.5\underline\sigma \mu\|\Delta^2v\|_2^2\nonumber \\
  &\leq& C\int(|[\Delta,\sigma]\nabla^2v||\nabla^4v|+|\nabla\sigma||\nabla^3v||\nabla^4v|)dxdydz+ \|\Phi_1\|_{H^2}\|v\|_{H^4}\nonumber\\
  &\leq& C\int(|[\Delta,\sigma]\nabla^2v||\nabla^4v|+|\nabla\sigma||\nabla^3v||\nabla^4v|)dxdydz\nonumber\\
  &&+ (\nu C_*\|\partial_z\sigma\|_{H^2}\|\partial_zv\|_{H^2}+C)\|v\|_{H^4},\label{EQ18}
\end{eqnarray}
where $C_*$ is an absolute positive constant.
Applying \eqref{COME} and by assumption \eqref{ASSUM-APRM}, it follows from the H\"older, Sobolev, and Gagliardo-Nirenberg inequalities that
\begin{eqnarray}
  &&\int|[\Delta,\sigma]\nabla^2v||\nabla^4v|dxdydz\nonumber\\
  &\leq& \|[\Delta,\sigma]\nabla^2v\|_2\|v\|_{H^4} \leq C\left(\|\nabla\sigma\|_6\|\nabla^2v\|_{W^{1,3}}+\|\sigma\|_{H^2}\|\nabla^2v\|_\infty\right)\|v\|_{H^4}\nonumber\\
  &\leq& C\|\sigma\|_{H^2}\|v\|_{H^3}^\frac12\|v\|_{H^4}^\frac12\|v\|_{H^4}\leq C\|v\|_{H^4}^\frac32 \label{EQ19}
\end{eqnarray}
and
\begin{eqnarray}
  \int|\nabla\sigma||\nabla^3v||\nabla^4v|dxdydz
  &\leq&\|\nabla\sigma\|_6\|\nabla^3v\|_3\|\nabla^4v\|_2\nonumber\\
  &\leq& C\|\sigma\|_{H^2}\|v\|_{H^3}^\frac12\|v\|_{H^4}^\frac32 \leq C\|v\|_{H^4}^\frac32. \label{EQ20}
\end{eqnarray}
Substituting (\ref{EQ19}) and (\ref{EQ20}) into (\ref{EQ18}) and adding the resultant with (\ref{EQ16}), it follows
\begin{equation}
  \frac{d}{dt}\|(v,\nabla\Delta v)\|_2^2+\underline\sigma \mu\|\Delta^2v\|_2^2
  \leq C(1+\|v\|_{H^4}^\frac32)+\nu C_*\|\partial_z\sigma\|_{H^2}\|v\|_{H^4}\|\partial_zv\|_{H^2},\label{EQ34}
\end{equation}
proving the first conclusion.


Next, we carry out the energy inequality for $\sigma$. Multiplying \eqref{EQs'} with $\sigma$
and integrating over $\mathcal O$, it follows from integrating
by parts, the H\"older and Sobolev inequalities, (\ref{ASSUM-APRM}), and Proposition \ref{PROP-ESTPhis} that
\begin{eqnarray}
  \frac12\frac{d}{dt}\|\sigma\|_2^2+\epsilon\|\nabla_h\sigma\|_2^2
  &=&\int(\nu\sigma\partial_z^2\sigma-v\cdot\nabla_h\sigma+\Phi_2)\sigma dxdydz\nonumber \\
  &\leq& C\|\sigma\|_\infty(\|\sigma\|_2\|\partial_z^2\sigma\|_2+\|v\|_2\|\nabla_h\sigma\|_2+\|\Phi_2\|_1)\nonumber\\
  &\leq& C(\|\sigma\|_{H^2}^3+\|v\|_2\|\sigma\|_{H^2}^2+\|\sigma\|_{H^2}\|\Phi_2\|_1)\leq C. \label{EQ23}
\end{eqnarray}
Applying $\Delta$ to \eqref{EQs'}, multiplying the resultant with $\Delta\sigma$, using the boundary conditions \eqref{BC1}, \eqref{BC2}, and \eqref{BC3},
it follows from integrating by parts that
\begin{eqnarray*}
  &&\frac12\frac{d}{dt}\|\Delta\sigma\|_2^2+\epsilon\|\nabla_h\Delta\sigma\|_2^2 \nonumber\\
  &=&\nu\int(\sigma\Delta\partial_z^2\sigma+[\Delta,\sigma]\partial_z^2\sigma)\Delta\sigma dxdydz+\int\Delta(\Phi_2-v\cdot\nabla_h\sigma)\Delta\sigma dxdydz\nonumber\\
  &=&-\nu\int(\sigma|\Delta\partial_z\sigma|^2+\partial_z\sigma\Delta\partial_z\sigma\Delta\sigma)dxdydz\nonumber\\
  &&-\nu\int([\Delta,\partial_z\sigma]\partial_z\sigma\Delta\sigma+
  [\Delta,\sigma]\partial_z\sigma\Delta\partial_z\sigma) dxdydz\nonumber\\
  &&+\int\Delta(\Phi_2-v\cdot\nabla_h\sigma)\Delta\sigma dxdydz
\end{eqnarray*}
and thus
\begin{eqnarray}
  &&\frac12\frac{d}{dt}\|\Delta\sigma\|_2^2+\epsilon\|\nabla_h\Delta\sigma\|_2^2+\nu\int\sigma|\Delta\partial_z\sigma|^2dxdydz \nonumber\\
  &=&-\nu\int \partial_z\sigma\Delta\partial_z\sigma\Delta\sigma dxdydz
 -\nu\int([\Delta,\partial_z\sigma]\partial_z\sigma\Delta\sigma+
  [\Delta,\sigma]\partial_z\sigma\Delta\partial_z\sigma) dxdydz\nonumber\\
  &&+\int\Delta(\Phi_2-v\cdot\nabla_h\sigma)\Delta\sigma dxdydz.\label{EQ25}
\end{eqnarray}
Using \eqref{ASSUM-APRM}, it follows from the H\"older, Sobolev, and Gagliardo-Nirenberg inequalities that
\begin{eqnarray}
\left|\int\partial_z\sigma\Delta\partial_z\sigma\Delta\sigma dxdydz\right|\leq \|\partial_z\sigma\|_\infty\|\Delta\partial_z\sigma\|_2\|\Delta
\sigma\|_2\nonumber\\
\leq C\|\partial_z\sigma\|_{H^1}^\frac12\|\partial_z\sigma\|_{H^2}^\frac32\|\sigma\|_{H^2}\leq C\|\partial_z\sigma\|_{H^2}^\frac32.\label{EQ26}
\end{eqnarray}
Applying \eqref{COME} and recalling \eqref{ASSUM-APRM}, it follows from the H\"older, Sobolev, and Gagliardo-Nirenberg inequalities that
\begin{eqnarray}
  \left|\int[\Delta,\partial_z\sigma]\partial_z\sigma\Delta\sigma dxdydz\right|
  &\leq&\|[\Delta,\partial_z\sigma]\partial_z\sigma\|_2\|\Delta\sigma\|_2\nonumber\\
  &\leq& C(\|\nabla\partial_z\sigma\|_3\|\partial_z\sigma\|_{W^{1,6}}
  +\|\partial_z\sigma\|_{H^2}\|\partial_z\sigma\|_\infty)\|\sigma\|_{H^2}\nonumber\\
  &\leq& C\|\partial_z\sigma\|_{H^1}^\frac12\|\partial_z\sigma\|_{H^2}^\frac32\|\sigma\|_{H^2}\leq C\|\partial_z\sigma\|_{H^2}^\frac32
  \label{EQ27}\\
  \left|\int[\Delta,\sigma]\partial_z\sigma\Delta\partial_z\sigma dxdydz\right|&\leq&\|[\Delta,\sigma]\partial_z\sigma\|_2\|\Delta\partial_z
  \sigma\|_2\nonumber\\
  &\leq&C(\|\nabla\sigma\|_6\|\partial_z\sigma\|_{W^{1,3}}+\|\sigma\|_{H^2}\|\partial_z\sigma\|_\infty)\|\partial_z\sigma\|_{H^2}\nonumber\\
  &\leq&C\|\sigma\|_{H^2}\|\partial_z\sigma\|_{H^1}^\frac12\|\partial_z\sigma\|_{H^2}^\frac32\leq C\|\partial_z\sigma\|_{H^2}^\frac32,
  \label{EQ28}
\end{eqnarray}
and
\begin{eqnarray}
  &&-\int\Delta(v\cdot\nabla_h\sigma)\Delta\sigma dxdydz\nonumber\\
  &=&-\int(v\cdot\nabla_h\Delta\sigma\Delta\sigma+[\Delta,v\cdot\nabla_h]\sigma\Delta\sigma)dxdydz\nonumber\\
  &\leq&\frac12\int\text{div}_hv|\Delta\sigma|^2dxdydz +\|[\Delta,v\cdot\nabla_h]\sigma\|_2\|\Delta\sigma\|_2\nonumber\\
  &\leq&C\|v\|_{H^3}\|\sigma\|_{H^2}^2+C(\|\nabla v\|_\infty\|\nabla_h\sigma\|_{H^1}+\|v\|_{W^{2,6}}\|\nabla_h\sigma\|_3)
  \|\Delta\sigma\|_2\nonumber\\
  &\leq&C\|v\|_{H^3}\|\sigma\|_{H^2}^2\leq C.\label{EQ29}
\end{eqnarray}
By Proposition \ref{PROP-ESTPhis} and recalling \eqref{ASSUM-APRM}, it follows
\begin{equation}
  \int\Delta\Phi_2\Delta\sigma dxdydz\leq C\|\Phi_2\|_{H^2}\|\sigma\|_{H^2}\leq C\left(1+\|\partial_z\sigma\|_{H^2}^\frac32\right). \label{EQ30}
\end{equation}
Substituting (\ref{EQ26})--(\ref{EQ30}) into \eqref{EQ25} and adding the resultant with \eqref{EQ23}, it follows from
\eqref{ASSUM-LBsp} and the Young inequality that
\begin{equation}
  \frac{d}{dt}\|(\sigma,\Delta\sigma)\|_2^2+\epsilon\|(\nabla_h\sigma,\nabla_h\Delta\sigma)\|_2^2
  +\nu\underline\sigma\|\Delta\partial_z\sigma\|_2^2\leq C\left(1+\|\partial_z\sigma\|_{H^2}^\frac32\right), \label{EQ31}
\end{equation}
proving the second conclusion.

Finally, we carry out the energy inequality for $p$. Multiplying \eqref{EQp} with $p$ and integrating over $\mathbb T^2$, it following from
Proposition \ref{PROP-ESTPhis} and \eqref{ASSUM-APRM} that
\begin{eqnarray}
  &&\frac12\frac{d}{dt}\|p\|_2^2+\epsilon\|\nabla_hp\|_2^2= \int(\Phi_3-\bar v\cdot\nabla_hp)pdxdy\nonumber
  \\
  &\leq& \|\Phi_3\|_1\|p\|_\infty+\|\bar v\|_2\|\nabla_hp\|_2\|p\|_\infty\leq C(\|p\|_{H^2}+\|v\|_2\|p\|_{H^2}^2)\leq C.
  \label{EQ32}
\end{eqnarray}
Applying $\nabla_h\Delta_h$ to \eqref{EQp}, multiplying the resultant with $\nabla_h\Delta_hp$, and integrating over $\mathbb T^2$, one
obtains by integrating by parts, \eqref{COME}, Proposition \ref{PROP-ESTPhis}, and \eqref{ASSUM-APRM} that
\begin{eqnarray*}
  &&\frac12\frac{d}{dt}\|\nabla_h\Delta_hp\|_2^2+\epsilon\|\nabla_h^2\Delta_hp\|_2^2\nonumber\\
  &=&\int\Big(\nabla_h\Delta_h\Phi_3-(\bar v\cdot\nabla_h)\nabla_h\Delta_hp-[\nabla_h\Delta_h,\bar v\cdot\nabla_h]p\Big)\cdot\nabla_h\Delta_hpdxdy \nonumber\\
  &\leq&\|\Phi_3\|_{H^3}\|p\|_{H^3}+\frac12\int\text{div}_h\bar v|\nabla_h\Delta_hp|^2dxdy+\|[\nabla_h\Delta_h,\bar v\cdot\nabla_h]p\|_2\|\nabla_h\Delta_h
  p\|_2\nonumber\\
  &\leq&C(1+\|v\|_{H^4})\|p\|_{H^3}+C\|\nabla v\|_\infty\|p\|_{H^3}^2\nonumber\\
  &&+C(\|\nabla_h\bar v\|_\infty\|\nabla_hp\|_{H^2}+\|\bar v\|_{H^3}\|\nabla_hp\|_\infty)\|p\|_{H^3}\nonumber\\
  &\leq&C(1+\|v\|_{H^4})\|p\|_{H^3}+C\|v\|_{H^3}\|p\|_{H^3}^2\leq C(1+\|v\|_{H^4}).
\end{eqnarray*}
Summing this with \eqref{EQ32} leads to
\begin{equation}
  \frac{d}{dt}\|(p,\nabla_h\Delta_hp)\|_2^2+\epsilon\|(\nabla_hp,\nabla_h^2\Delta_hp)\|_2^2
  \leq C(1+\|v\|_{H^4}),\label{EQ33}
\end{equation}
proving the third conclusion.
\end{proof}

\end{document}
